\section{Mathematical Appendix}
\label{appendix:theory}

We collect here the formal conditions under which our estimator is consistent and asymptotically normal. These results are adaptations of classical results from the theory of just-identified generalized method of moments (GMM) estimators, tailored to the structure of our estimating equations.

\subsection{From Marginal Log-Likelihood to MLE Estimating Equations}
\label{apx:mle-estimating-equations}
To derive the estimating equations from the marginal log-likelihood, we proceed as follows:

\begin{enumerate}
    \item \textbf{Start with the marginal log-likelihood:}
    \[
    \ell(\bt; \y_i) = \log \left( \int_{\mathcal Z} f_\bt(\y_i, \z)\,d\z \right),
    \]
    where \( f_\bt(\y_i, \z) = f_\bt(\y_i \mid \z) f(\z) \).

    \item \textbf{Differentiate under the integral sign:}  
    The derivative of the log-likelihood w.r.t. \( \bt \) is
    \[
    \nabla_\bt \ell(\bt; \y_i) = \frac{ \int_{\mathcal Z} \nabla_\bt f_\bt(\y_i, \z)\, d\z }{ \int_{\mathcal Z} f_\bt(\y_i, \z)\, d\z }.
    \]

    \item \textbf{Rewrite the numerator using the score function:}
    \[
    \nabla_\bt f_\bt(\y_i, \z) = f_\bt(\y_i, \z) \nabla_\bt \log f_\bt(\y_i, \z).
    \]
    Substituting gives:
    \[
    \nabla_\bt \ell(\bt; \y_i)
    = \frac{ \int f_\bt(\y_i, \z) \nabla_\bt \log f_\bt(\y_i, \z)\, d\z }{ f_\bt(\y_i) }
    = \int \nabla_\bt \log f_\bt(\y_i, \z) \cdot \frac{f_\bt(\y_i, \z)}{f_\bt(\y_i)}\, d\z.
    \]

    \item \textbf{Recognize the posterior:}  
    The ratio \( \frac{f_\bt(\y_i, \z)}{f_\bt(\y_i)} \) is exactly the posterior density:
    \[
    f_\bt(\z \mid \y_i) := \frac{f_\bt(\y_i, \z)}{f_\bt(\y_i)}.
    \]

    \item \textbf{Final expression:}
    \[
    \nabla_\bt \ell(\bt; \y_i)
    = \mathbb{E}_{\Z_i \sim f_\bt(\z \mid \y_i)} \left[ \nabla_\bt \log f_\bt(\y_i, \Z_i) \right].
    \]

    \item \textbf{Stacking over all \( n \) observations:}
    \[
    \frac{1}{n} \sum_{i=1}^n \mathbb{E}_{\Z_i \sim f_\bt(\z \mid \y_i)} \left[ \nabla_\bt \log f_\bt(\y_i, \Z_i) \right] = 0,
    \]
    which gives the MLE estimating equations as in Equation~\eqref{eqn:MLE-estimating-equations}.
\end{enumerate}





\subsection{Derivation of the Jacobian}
\label{math:jacobian}
The Jacobian takes the form,
%
\begin{align}
A_{q,c}(\bt_0)\;&=\;-\nabla_\bt \Psi(\bt)\big|_{\bt=\bt_0}\\
&=-\E_{\bt_0}\left[\nabla_\bt \psi_{q,c}(\bt;\Y)\right]_{\bt=\bto} + \E_{\bt_0}\left[\nabla_\bt \psi_{q,c}(\bt;\Y)\right]_{\bt=\bto} \\
&\quad\, - \E_{\bt}\left[\psi_{q,c}(\bt;\Y)s(\bt;\Y)^\top\right]_{\bt=\bto}\\
&=-\E_{\bt_0}\left[\psi_{q,c}(\bt_0;\Y)\,s(\bt_0;\Y)^\top\right]\\
&=-\E_{\bt_0}\left[\big(m_q(\Y;\bt_0)-c\,b_q(\Y;\bt_0)\big)\,s(\bt_0;\Y)^\top\right],
\end{align}
%
Where we used the fact, that, for any measurable function $g(\bt,\Y)$,
\[
\E_{f_{\Y;\,\bt}}[g(\bt,\Y)] \;=\; \int g(\bt,y)\, f_{\Y;\,\bt}(y)\,dy.
\]
Differentiating w.r.t.\ $\bt$ under the integral sign gives
\begin{align}
\nabla_\bt \E_{f_{\Y;\,\bt}}[g(\bt,\Y)]
&= \int \nabla_\bt g(\bt,y)\, f_{\Y;\,\bt}(y)\,dy \;+\; \int g(\bt,y)\,\nabla_\bt f_{\Y;\,\bt}(y)\,dy\\
&= \E_{f_{\Y;\,\bt}}[\nabla_\bt g(\bt,\Y)] \;+\; \E_{f_{\Y;\,\bt}}\!\Big[g(\bt,\Y)\,\frac{\nabla_\bt f_{\Y;\,\bt}(\Y)}{f_{\Y;\,\bt}(\Y)}\Big]\\
&= \E_{f_{\Y;\,\bt}}[\nabla_\bt g(\bt,\Y)] \;+\; \E_{f_{\Y;\,\bt}}[g(\bt,\Y)\,s(\bt;\Y)^\top],
\end{align}
where $s(\bt;\Y)=\nabla_\bt \log f_{\Y;\,\bt}(\Y)$ is the score function.


\section{OTHER STUFF}

\subsection{MLE compatibility intuition}

A useful compatibility check comes from the exact-score case. Under the true model, the
likelihood score equation implies that at $\bt=\bt_0$ the two posterior expectations in
\eqref{eq:score_two_terms} have the same population target:
\begin{equation}
\mathbb E_{f_{\Y;\bt_0}}
\!\left[
\mathbb{E}_{f_{\Z\mid\Y;\bt_0}(\cdot\mid \Y)}\!\big[D(\Z;\bt_0)^\top T(\Y)\big]
\right]
=
\mathbb E_{f_{\Y;\bt_0}}
\!\left[
\mathbb{E}_{f_{\Z\mid\Y;\bt_0}(\cdot\mid \Y)}\!\big[D(\Z;\bt_0)^\top \mu(\Z;\bt_0)\big]
\right].
\label{eq:mle_pop_matching_bt0}
\end{equation}
If $q_{\Z\mid\Y}=f_{\Z\mid\Y;\bt_0}$, then
$\psi_q(\bt_0;\Y)=\mathbb{E}_{f_{\Z\mid\Y;\bt_0}(\cdot\mid \Y)}[D(\Z;\bt_0)^\top T(\Y)]$, so
the model-based centering term in \eqref{eq:zq_matching} evaluated at $\bt_0$ coincides with
the population target appearing on the left-hand side of \eqref{eq:mle_pop_matching_bt0}, and
hence also with the population target of the MLE centering term on the right-hand side.

For an arbitrary surrogate $q_{\Z\mid\Y}$, we do not expect such an identity to hold pointwise.
Instead, the centering in \eqref{eq:Psi_nq} enforces the population condition
$\Psi_q(\bt_0)=0$ by construction, which is the basis for decoder consistency under surrogate
misspecification.








